\documentclass[11pt]{article}

\usepackage[T1]{fontenc}
\usepackage[utf8]{inputenc}
\usepackage{amsmath,amssymb}
\usepackage[margin=1in]{geometry}
\usepackage{xcolor}
\usepackage{graphicx}
\usepackage{booktabs}
\usepackage{longtable}
\usepackage[round,authoryear]{natbib}
\usepackage[hidelinks]{hyperref}
\graphicspath{{paper/figures/}}
\AddToHook{env/figure/begin}{\centering}
\AddToHook{env/table/begin}{\centering}
\setlength{\LTleft}{\fill}
\setlength{\LTright}{\fill}

\definecolor{boxrule}{HTML}{2C6EAB}
\newcommand{\fit}{\phi}
\newcommand{\data}{\mathcal{D}}
\newcommand{\resolved}{\rho_{\alpha}}
\newcommand{\rank}{\mathrm{rank}}
\newcommand{\note}[1]{\textcolor{boxrule}{\small[#1]}}

\title{How well do deep mutational scans resolve top-ranked variant batches?\\
\large Ranking uncertainty in protein engineering and machine-learning benchmarks}
\author{[TODO: authors and affiliations]}
\date{Working draft \textbar\ September 2026}

\begin{document}
\maketitle

\noindent\note{Working draft. Missing experimental records, model checks,
and publication metadata are identified in the text and supplement.}

\begin{abstract}
Protein engineers use measurements of protein fitness, performance in a
specified experimental test, to choose variants for further study.
Machine-learning researchers use the same measurements to train and evaluate
predictors that help select candidate batches.
Deep mutational scans measure many variants together, and sorting their
reported scores gives a convenient ordering.
But that ordering depends on how fitness is calculated from sequencing
counts, and an exact list does not show the uncertainty left by finite counts.
In the present work, we examine how these choices and uncertainties affect
leading ranks in three protein systems measured through growth or fluorescence.
We use fitted count models to generate joint sets of fitness scores and
measure rank intervals, candidate inclusion frequencies, and overlap between
selected batches.
The fitted and reported score summaries select different batches, and
repeated sampled selections can also replace a substantial fraction of a
batch within a fixed analysis.
The extent of turnover varies across landscapes and selection budgets,
so an exact list of leading variants can conceal uncertainty
about which candidates to select.
\end{abstract}

\section{Introduction}
\label{sec:introduction}

Protein engineers change amino-acid sequences to improve a chosen function.
For example, they may seek an enzyme that carries out a chemical reaction
more effectively or a protein that binds a particular partner
\citep{yang2019}. Performance in the experimental test chosen for that
purpose is often called \emph{protein fitness}. Improving that performance
requires finding useful sequence changes without testing every possible
combination. Allowing each of four positions to take any of the twenty
standard amino acids already gives 160,000 combinations. In directed
evolution, measurements guide successive rounds of sequence changes,
but a change that helps in one sequence background need not help in
another. This dependence, called epistasis, makes combinations of changes
important to the search \citep{yang2019,johnston2024}. Testing combinations
therefore competes with the need to measure each one well.

Machine learning helps organize this search. A predictor trained on measured
sequences and fitness scores can propose a batch of untested variants for
the next experiment. The choice of training data and the ability to predict
beyond it both affect this process \citep{yang2019,wittmann2021}.
The same sequence--score pairs are used to evaluate predictors: an analyst
asks whether a model recovers the measured ordering or identifies the
leading candidates. Benchmarks such as FLIP and ProteinGym collect
experimental data for these comparisons \citep{dallago2021,notin2023}.
The interpretation of a prediction therefore depends on what its fitness
label measures and how well the experiment determines that label.

The immediate output of this workflow is often a batch, not an exact
ordering of every possible sequence. An engineer with resources to test
a hundred variants needs a candidate group; the order of two variants
within that group may not affect the experiment. A researcher evaluating
the predictor may instead measure how closely it reproduces the entire
reported ordering. These objectives are related but not equivalent.
Separating clearly different variants can give a useful overall
correlation while leaving the choice among closely competing leaders
unresolved. The relevant experimental resolution depends on the decision
the predictions will support.

Training and evaluation use these measurements in different ways. During
training, experimental scores determine which sequence changes the model
learns to associate with improved performance. During evaluation, the
reported scores become the reference against which predictions are judged.
Uncertainty in that reference is distinct from uncertainty in the predictor.
Even a predictor that exactly reproduced a released table would inherit
the table's decisions about fitness construction and the resolution of its
leading ranks. This does not make the labels unusable. It makes their
experimental interpretation part of the prediction problem, especially
when the intended output is a short list of candidates for testing.

Fitness is defined by an assay, not by sequence alone. An assay is a
specified experimental test, and different tests can summarize different
properties of the same protein. Binding or enzyme activity may be coupled
to the growth of cells carrying a variant; activity may instead produce
a fluorescent signal used to select cells or droplets. These measurements
do not all estimate one universal quantity. For example, growth supported
by an enzyme can depend on both its activity and the amount of functional
protein in the cell. A fitness score from that experiment is not necessarily
a direct measurement of the purified enzyme's reaction rate.

Deep mutational scanning makes such measurements for many sequence variants
in parallel \citep{fowler2014}. The experiment starts with a \emph{library}
of DNA sequences encoding different protein variants. After a pooled
selection, researchers sequence the DNA to identify which variants became
more common. Sequencing before selection gives \emph{input counts};
sequencing the selected population gives \emph{output counts}. A read is
an observation of a DNA fragment assigned to a variant, not a direct
measurement of that protein's fitness. A common score begins with
\emph{enrichment}, the variant's fraction of output reads divided by its
fraction of input reads. Using both populations accounts for unequal
starting abundances, while using fractions accounts for unequal total
numbers of reads, or sequencing depths.

The physical readout and its statistical summary are also different
choices. A growth assay can be summarized by enrichment at an endpoint
or by a growth parameter fitted across observations. A fluorescence
screen can be summarized by passage through one threshold or by a
signal distribution fitted from several thresholds. These choices use
the observations to answer related questions, but they need not induce
the same ordering. Before asking whether a ranking is precise, we need
to specify which assay summary is being ranked.

The enrichment ratio alone hides how much data supports it. In a
hypothetical experiment with equal input and output depths, 10 output reads
for 1 input read and 10,000 output reads for 1,000 input reads both give
a ratio of ten. One extra input read halves the first ratio and barely
changes the second. A released score can therefore have different precision
for well-observed and sparsely observed variants. Combining observations
also requires care: a new biological population, another sequencing sample
from the same population, a later timepoint, and a different fluorescence
threshold are not interchangeable replicates.

The observation structure determines what additional information the
counts provide. Biological replicates can reveal variation between
populations, timepoints describe changes during selection, and multiple
gates probe different parts of a signal distribution. Simply adding all
these counts would discard those distinctions. A statistical model can
retain them while linking the counts to variant-level performance.
Uncertainty in the resulting scores then reflects the information
available under that model, rather than an identical error assigned to
every variant.

A released \emph{fitness landscape} usually contains sequences and one
reported score per variant. Sorting those scores gives the \emph{reported
ordering}. Taking its leading hundred candidates is a natural way to
propose a validation batch or assess a predictor's retrieval of strong
variants. But two questions remain. First, would another calculation of
fitness from the same observations select the same candidates? Scores
depend on choices about normalization, reference sequences, and how
observations are combined. Second, how much would the proposed batch vary
under the uncertainty in one specified analysis? An exact sorted list
answers neither question.

Methods for estimating fitness already address measurement uncertainty.
Enrich2 and DiMSum estimate score errors from counts and repeated
measurements \citep{rubin2017,faure2020}; model-based enrichment estimation
shares information between related sequences \citep{busia2023}. Analyses
of experimental design also study the effects of read depth and sampling
choices on precision \citep{matuszewski2016}. Our question is what the
remaining uncertainty means for ranks and batches. Separate score intervals
do not answer it: shared shifts of all scores leave ranks unchanged,
whereas relative changes between close competitors can exchange their
positions. Two variants can repeatedly exchange ranks and both stay in
a batch; a smaller movement across its cutoff can replace a member.

This leaves an empirical question beyond whether the measurements are
noisy: does their uncertainty mainly concern exchanges within a robust
group, or which variants belong to that group? The answer depends on
the competing scores, how they vary together, and the selection budget.
Increasing a budget also moves the cutoff to a different part of the
landscape. A precision estimate for one score, or a single correlation
over the landscape, cannot determine the answer at every cutoff.
We therefore assess the sampled selections themselves.

In the present work, we analyze sequencing counts from three protein
systems: PDZ-mediated binding, tryptophan production by TrpB, and DNA
cleavage by NucB \citep{zarin2024,johnston2024,thomas2025}. These landscapes
are also used in the FLIP2 protein-fitness benchmark \citep{didi2026}.
Their distinct experimental readouts require corresponding count models.
We use these models to generate joint sets of fitness scores, rank each set against
all eligible competitors, and select its leading batch. We examine
(1) agreement with reported leading ranks, (2) uncertainty in ranks and
batch membership within an analysis, and (3) the count support and score
differences associated with plausible replacements. Reported top-100 and
top-1,000 cohorts and batch budgets from 10 to 1,000 keep the analysis
focused on leading candidates.

The fitted and reported score summaries select different batches, and
agreement within a fitted analysis does not imply agreement with the
reported list. Within an analysis, uncertain individual ranks can coexist
with a largely stable batch, and a larger batch budget does not always
increase the shared fraction of candidates. Lower input counts are
associated with wider score intervals in all analyzed leading-candidate
cohorts, but interval width alone does not identify which candidates are
replaceable. These results concern the experimental labels before a
predictor is introduced.

\section{Results}
\label{sec:results}

\subsection{Three experimental readouts of protein fitness}
\label{sec:overview}

The three systems connect protein function to sequencing counts in
different ways. The PDZ assays couple binding between PSD95-PDZ3 and
its partner CRIPT to yeast growth; their reported scores summarize
binding-dependent changes in abundance \citep{zarin2024}. In the TrpB
assays, introduced enzyme variants support bacterial growth by producing
tryptophan, and sequencing follows their abundance over time
\citep{johnston2024}. NucB activity is measured through fluorescence
sorting: a \emph{gate} collects a population above a specified signal
threshold, and sequencing measures which variants enter it
\citep{thomas2025}. We analyze five PDZ landscapes and four NucB
engineering generations. The TrpB data contain nine libraries varying
three amino-acid positions and one varying four positions.
These are different candidate collections, not replicate experiments of
one landscape. Studying both growth and sorting assays lets us examine
the ranking questions across distinct ways of turning protein performance
into sequencing counts, without equating their fitness units.
Supplementary Section~\ref{app:experimental-records}
explains the assays and their reported labels in detail.

We fit growth models to PDZ and TrpB counts and signal-distribution models
to NucB counts. Each model accounts for a variant's input abundance when
describing its output counts. A posterior distribution represents parameter
uncertainty conditional on the observed counts and the model.
We use three kinds of samples to ask how ranking conclusions vary.
\emph{Parameter resampling} draws from this posterior to calculate
fitness. \emph{Optimized-fit
count resampling} generates alternative count observations from a single
optimized fit and recalculates count-based fitness. \emph{Posterior count
resampling} also generates alternative counts, but varies the parameters
over their posterior distribution. Both count-based recipes retain the
recorded sequencing depths and do not
refit the model for each draw.

Parameter samples ask how uncertain the fitted performance parameters
are. Count samples ask how the count-derived scores would vary across
alternative sequencing outcomes at the recorded depths. We keep those
questions separate rather than interpreting the recipes as interchangeable
estimates of one source of noise. None simulates every stage of a new
physical experiment. Each draw keeps all candidate
scores together, including their shared variation. Its ranks therefore
compare each candidate with every finite, eligible competitor in that
draw. We retain 158 analyses: 49 parameter, 60 optimized-fit count, and
49 posterior count runs. A \emph{run} is one landscape, model specification,
and sampling recipe; several runs can use the same observations.
Materials and Methods gives the sampling formulas and inventory.

We use the same three running examples throughout: PDZ trans-1 and
TrpB libB under Exponential--Exponential, and NucB G3 under Exponential.
The examples illustrate different relationships between score precision,
rank uncertainty, and batch turnover; they were not chosen as median cases.
The growth-model names distinguish rate assumptions
and growth functions; the NucB name distinguishes signal-location
assumptions. Methods Section~\ref{methods:examples} gives the selection
rule and its limitations. All model specifications remain in the overview
panels and supplement.

\subsection{Reported and fitted top-100 lists differ}
\label{sec:reported}

The fitted and reported scores select different top-100 batches. To make
this comparison, we first take each candidate's median sampled score and
rank these medians. We call the resulting fixed list the
\emph{median-score reference}. This gives a proposed batch from the
fitted analysis to compare with the batch obtained by sorting reported
scores. Under parameter resampling, its top 100
shares 52 candidates (52\%) with the reported batch for PDZ trans-1, 33 (33\%)
for TrpB libB, and 68 (68\%) for NucB G3
(Table~\ref{tab:reference-comparisons}). These fixed-list intersections
compare two summaries of the same experimental data, before examining
turnover among individual sampled batches.

For a hundred-candidate validation budget, the two summaries therefore
agree on about half the candidates in PDZ trans-1, one third in
TrpB libB, and two thirds in NucB G3. Within each example, both summaries
use the same source experiment. The disagreement is
not between competing sequence predictors; it arises before a predictor
is introduced, when the experimental observations are converted into
the list against which predictions might be judged.

The scores need not estimate the same assay summary. For NucB, the
reported comparator uses enrichment at the highest selected gate, while
the fitted parameter score uses counts from all available gates. For PDZ,
our count-based calculation approximates the original DiMSum score.
Disagreement does not, by itself, show that the reported ordering is wrong.

Changing the units of a fixed score list alone would not explain this
result: a common strictly increasing transformation preserves its ordering. Here the
calculations also differ in how they combine observations and which
performance parameter or count summary they estimate. In TrpB, for
example, the reported calculation normalizes count ratios relative to
stop-codon variants and an observation's maximum before averaging;
the parameter score orders growth parameters fitted to the time-indexed
counts. The fixed-list comparison establishes the consequence for
selection without assuming that the two scores are interchangeable.

We next ask where each reported rank falls within that candidate's sampled
rank distribution. This adds information absent from a fixed-list
intersection: do reported leaders fall near the middle of their sampled
ranks, or do most sampled rankings move them in one direction? For a
candidate, its \emph{reported-rank sample
percentile} is the fraction of finite sampled ranks better than the
reported rank, giving ties half weight. A value near zero means that the
reported rank is better than most sampled ranks; a value near one means
it is worse. The curves in Fig.~\ref{fig:reported}a--c summarize these
percentiles over the reported top 100. At each horizontal value, the
vertical value is the fraction of candidates with a percentile at or
below it. A steep increase near zero therefore identifies reported
leaders that usually move down in the fitted analysis.

That increase is evident in all three examples
(Fig.~\ref{fig:reported}a--c). Under parameter resampling, 65 of the
reported top 100 TrpB candidates (65\%) have percentiles at or below
0.1, compared with 55 (55\%) for NucB G3 and 42 (42\%) for PDZ trans-1.
The concentration near zero therefore involves most of the reported
TrpB leaders, not a small subset of the cohort.
Both count-based TrpB curves have the same broad shape: 63 of 100
candidates under optimized-fit count resampling and 65 of 100 under
posterior count resampling have percentiles at or below 0.1.
These fractions describe where reported ranks lie within the sampled
distributions; they do not measure the number of rank positions moved.

PDZ trans-1 shows a stronger difference between sampling recipes
(Fig.~\ref{fig:reported}a). Its optimized-fit count curve places 67 of
100 percentiles at or below 0.1, compared with 42 for parameter and
37 for posterior count resampling. The curves also rise at different
rates through the interior. They retain the direction and shape of
these differences before either is reduced to an area or correlation.

The middle row of Fig.~\ref{fig:reported} summarizes each curve by the
absolute area between it and the uniform diagonal. Under parameter
resampling, the areas are 0.210 for PDZ trans-1, 0.314 for TrpB libB, and
0.232 for NucB G3. Area provides one value per analysis for the overview,
but smaller area is not automatically better rank agreement. Exact
recovery of every reported rank would put every percentile at 0.5,
not produce a uniform distribution. Area also discards the direction
and shape of a departure. Table~\ref{tab:reported-curve-statistics}
reports each analysis's area and correlation for both cohorts.
Supplementary Section~\ref{app:all-reported-curves} groups the full
curves by protein system and sampling scheme.

Relative order is a separate question. The bottom row of
Fig.~\ref{fig:reported} shows the Spearman correlation between reported
ranks and median sampled ranks within the reported top 100. Spearman
correlation measures whether candidates occur in a similar relative
order; it does not measure how far their ranks move or how variable
those ranks are. Across the three recipes, correlations are
0.190--0.319 for PDZ trans-1, 0.050--0.056 for TrpB libB, and 0.503--0.584
for NucB G3. The black outlines and arrows identify these exact fits
among the individual model points, rather than identifying a
median across models as the example.

Each small filled dot in the overview panels is one summary for a
landscape, model specification, and sampling scheme. It is not a posterior
draw, chain, variant, or experimental replicate. For example, TrpB libB
contributes three model dots per scheme, whereas NucB G3 contributes one.
Larger hollow symbols mark the median across models within each landscape
and scheme. Small fixed vertical offsets help distinguish nearby model
points; their horizontal metric values are unchanged. When only one model
is available, its dot lies inside the hollow median symbol.

The overview rows also show why one example cannot stand for an entire
protein system. Under Exponential--Exponential parameter resampling,
the reported-top-100 correlation is 0.943 for PDZ cis and 0.310 for
PDZ trans-1 (Fig.~\ref{fig:reported}g). In TrpB it is 0.904 for
libD but $-0.082$ for libC under the same specification (panel h).
Individual model specifications are retained because changing the model
can also change the result within a landscape. The hollow symbols summarize
models fitted to shared data, not independent experimental replications.

The cohort also matters. For NucB G3, expanding the reported cohort to
1,000 raises the correlations to 0.921--0.926; agreement over that larger
group hides disagreement within its top 100. TrpB libB also improves
over the larger cohort, but its correlations remain only 0.252--0.262
(Fig.~\ref{fig:reported1000}). A predictor evaluated over a larger
cohort would therefore face a different ranking comparison, even using
the same underlying experiment. The area and correlation rows address
complementary aspects of that comparison: departure of the percentile
distribution and preservation of relative order. Neither quantifies
how often an individual candidate is selected within the fitted analysis.

\begin{figure}[p]
\centering
\includegraphics[width=\linewidth]{reported_order.pdf}
\caption{Reported-rank percentiles, area, and relative order.
Panels (a)--(c) show cumulative distributions of reported-rank sample
percentiles for the reported top 100 in PDZ trans-1 and TrpB libB
(both Exponential--Exponential), and NucB G3 (Exponential).
Panels (d)--(f) show the absolute area between each curve and the
diagonal; panels (g)--(i) show reported-versus-median-sampled-rank
Spearman correlations. Each small filled dot summarizes one model
specification within a landscape and sampling scheme, not an individual
draw. Hollow blue circles, red squares, and green triangles mark model
medians for parameter, optimized-fit count, and posterior count resampling
(49, 60, and 49 analyses). Vertical offsets improve visibility only;
horizontal values are unchanged. Black outlines and arrows mark example
model points. The diagonal is descriptive, not a calibration target.
Complete curves appear in Supplementary Section~\ref{app:all-reported-curves}.}
\label{fig:reported}
\end{figure}

\paragraph{How far do the reported ranks move?}
The percentile curves describe where a reported rank falls within its
sampled distribution, but not whether the difference is a few positions
or hundreds. We therefore also calculate \emph{mean absolute rank
displacement}: for each candidate in the reported top 100, we take the
absolute difference between its reported rank and its median sampled
rank, then average over supported candidates. Table~\ref{tab:rank-displacement}
reports this quantity for every landscape/model specification, with the
three sampling schemes in separate columns. The ranks still refer to
the full eligible competitor pool, so their displacement can exceed
the size of the reported cohort.

Under parameter resampling, the mean absolute displacement is 427.6
positions for PDZ trans-1, 1,116.1 for TrpB libB, and 42.2 for NucB G3.
The difference persists across sampling schemes: the example ranges are
427.6--958.7, 1,116.1--1,600.9, and 42.2--44.2 positions, respectively.
TrpB libB's low correlation in Fig.~\ref{fig:reported}h therefore
accompanies large differences in rank, not only exchanges within its
reported top 100. NucB G3 also differs from its reported ordering, but
the absolute movements are much smaller in these analyses.

Displacement and percentile area distinguish these examples differently.
PDZ trans-1 has a smaller parameter-resampling area than NucB G3
(0.210 versus 0.232), yet a larger mean absolute displacement
(427.6 versus 42.2 positions). A percentile near an endpoint can
result from a small but consistent rank difference as well as a large
one. The area cannot distinguish their sizes. Correlation addresses
another aspect: whether candidates retain their relative order, which
can persist even when other candidates move ahead of the reported
cohort. The table supplies the rank-distance information absent from
both Figure~\ref{fig:reported} overview rows.

PDZ CRIPT-C also distinguishes correlation from displacement. Under
parameter resampling, its Exponential--Exponential and Gamma--Exponential
analyses have reported-top-100 correlations of 0.885 and 0.887,
but mean absolute displacements of 12.9 and 316.4 positions, respectively.
These similarly high correlations accompany different amounts of
movement in the full ranking.
Table~\ref{tab:rank-displacement} retains these model-specific
differences rather than averaging them into one landscape value.

These displacements compare reported ranks with the centers of sampled
rank distributions; they do not measure the spread of those distributions.
A value of zero would mean that all supported median sampled ranks equal
their reported ranks, not that the ranks are certain. Large movements
by a few candidates can also influence the mean, and different landscapes
have different numbers of eligible competitors. We therefore read the
table as an absolute rank-position comparison on each landscape, not
as a normalized error against biological truth. The following rank
intervals and batch summaries address the uncertainty and selection
consequences separately.

\clearpage
\input{paper/figures/rank_displacement.tex}

\subsection{Broad rank intervals can coexist with stable batches}
\label{sec:noise}
\label{sec:ranks}

Leading candidates can have uncertain positions while remaining in
the selected batch. We describe each candidate's rank distribution
by its central 90\% interval, from the 5th to the 95th sampled-rank
percentile. Among the reference top 100, the median interval width
under parameter resampling is 178.5 positions for PDZ trans-1, 281 for
TrpB libB, and 38 for NucB G3. These are widths of individual rank
distributions, not error bars on an average rank.

Figure~\ref{fig:ranks}a--c displays the first 20 candidates in each
median-score reference. Blue points show median sampled ranks;
black ticks show ranks of median scores. They can differ because
ranking before taking a median is not the same as taking medians
before ranking. The intervals can extend farther toward worse ranks
than toward better ranks, and a reference rank need not fall inside
its candidate's sampled-rank interval. The bottom row summarizes
widths for the reference top 100 in every analysis. For TrpB libB,
the median width increases from 281 under parameter resampling to
624.2 under optimized-fit count resampling and 859.5 under posterior
count resampling. Moving from parameter to count resampling changes
both the score construction and the uncertainty represented; the
difference is not a pure decomposition of experimental noise.

The NucB G3 intervals illustrate why their location matters as much as
their width. The candidate at reference rank 10 has a median sampled
rank of 21 and a central 90\% interval from rank 3 to rank 40
(Fig.~\ref{fig:ranks}c). Across the first 20 reference candidates,
median sampled ranks lie between 16 and 23, and the intervals overlap
extensively. Each candidate occupies a broad range of rank positions;
the marginal intervals do not specify how each pair is ordered jointly.
Yet all 20 remain in the sampled top 100
in every retained parameter draw. Uncertainty about their individual ranks is not
uncertainty about their membership in that larger batch.

The conclusion changes closer to the cutoff. Among the full reference
top 100, 78 PDZ trans-1 candidates, 84 TrpB libB candidates, and 18 NucB G3
candidates have a 95th-percentile rank beyond 100 under parameter
resampling. Thus the first 20 intervals do not describe every member
of the proposed batch. The lower panels summarize the broader group,
while the inclusion analysis below measures how often each candidate
actually crosses its boundary.

The landscape summaries in Fig.~\ref{fig:ranks}d--f put those examples
in context. After taking the median across models within each TrpB
landscape, the reference-top-100 rank width under parameter resampling
ranges from 12 positions for libD to 285.5 for libA. The small filled
points also reveal specifications
with much wider intervals than neighboring fits of the same landscape.
This variation is a reason to inspect models separately, not evidence
that all their intervals are equally adequate descriptions of the
experiment. Sampler summaries and the status of predictive checks are discussed
in Supplementary Section~\ref{app:fit-checks}.

\begin{figure}[p]
\centering
\includegraphics[width=\linewidth]{rank_uncertainty.pdf}
\caption{Rank uncertainty within leading batches. Panels (a)--(c)
show the first 20 candidates in the median-score references for PDZ
trans-1 and TrpB libB (both Exponential--Exponential), and NucB G3
(Exponential). Intervals span the 5th--95th sampled-rank
percentiles under parameter resampling; blue points mark median
sampled ranks and black ticks mark median-score reference ranks.
Horizontal scales differ between examples. Panels (d)--(f) show
the median interval width among each run's reference top 100.
Small filled dots each summarize one model specification within a
landscape and sampling scheme; hollow symbols mark medians across
those models. Vertical offsets are for visibility only. These dots
are not individual posterior draws. Black outlines and arrows mark
the exact parameter fits displayed above. Blue circles, red squares,
and green triangles denote parameter, optimized-fit count, and
posterior count resampling. Logarithmic horizontal axes report
rank positions; the dashed line marks a width of 100. Intervals
describe sample distributions, not validated coverage of a known
biological ordering.}
\label{fig:ranks}
\end{figure}

\paragraph{Batch membership.}
\label{sec:decisions}
A rank movement changes a batch only if it crosses the selection
cutoff. For a budget of $k$ candidates, we select the top $k$
separately in each joint score sample. A candidate's \emph{inclusion frequency} is the
fraction of retained samples in which it is selected. The
\emph{mean shared fraction}, denoted $\widehat C_k$, is the average
fraction of members shared by two distinct sampled size-$k$ batches.
For example, a value of 0.8 at $k=100$ means 80 shared members on
average, not an 80\% chance that the complete lists match.

These two summaries serve different decisions. Inclusion frequency
locates the candidates whose membership changes; shared fraction
describes how much of the batch changes overall. The same shared
fraction could arise from many low-frequency alternatives or a
smaller number of frequently exchanged candidates. A batch-level
number therefore needs a candidate-level view to identify its
uncertain members.

For the three running examples, sampled top-100 batches under
parameter resampling share approximately 60, 51, and 93 members
(60\%, 51\%, and 93\%), respectively. In PDZ and TrpB, uncertainty
therefore changes much of the membership, not only the internal order.
Across the three recipes, TrpB libB shares about 26--51 members on
average, compared with
87--93 for NucB G3 (Fig.~\ref{fig:selection}e,f). These are comparisons
within each fitted sampling distribution, not observed agreement
between new physical experiments.

The top row of Fig.~\ref{fig:selection} locates the candidates that
can cross the cutoff. It shows reference ranks 91--100 and the ten
outside candidates with the largest inclusion frequencies under
parameter resampling. These candidates can enter a sampled batch
despite lying outside the fixed median-score reference. We call an
outside candidate a \emph{challenger} when its top-100 inclusion
frequency is at least 0.01. Challengers need not belong to the
reference top 1,000.

For TrpB libB, the ten reference candidates near the cutoff have
inclusion frequencies of 28--43\%, while the ten displayed outside
candidates have frequencies of 30--34\% (Fig.~\ref{fig:selection}b).
The most frequently selected outside candidate exceeds five of those
ten reference members. The median-score ordering is therefore not
an ordering by inclusion frequency. In PDZ trans-1 the displayed ranges
are 29--43\% inside and 33--40\% outside (panel a); in NucB G3 they
are 54--93\% and 18--46\% (panel c). The NucB boundary separates
these displayed groups more clearly, but membership of some reference
candidates is still far from certain. A hard cutoff in a fixed list
does not create a corresponding certainty threshold.

Agreement within an analysis does not imply agreement with the
reported batch. A parameter-sampled NucB G3 batch shares 95.36
members on average with its median-score reference, but 68.14
with the reported batch. For TrpB libB, those numbers are 63.44
and 29.81, while two sampled batches share 50.96 members
(Table~\ref{tab:reference-comparisons}). NucB therefore mainly differs
from the reported list while retaining a comparatively stable sampled
batch; TrpB also has substantial variation within the fitted analysis.
Overlap with the reported batch alone would not distinguish those cases.

The PDZ row combines both effects. Its parameter samples share 70.59
members with the median-score reference, 47.51 with the reported batch,
and 59.80 with another sampled batch. The \emph{Reference} column
compares samples with one fixed summary of the fitted score distribution;
\emph{Sample pair} compares two variable selections. Neither is the
fixed intersection of 52 candidates. Keeping these columns separate
distinguishes disagreement between summaries from variation around them.

The distinction persists across sampling recipes. TrpB libB's fixed
intersection changes only from 33 to 34 candidates, while its sample-pair
overlap changes from 50.96 under parameter resampling to 32.12 and
25.87 under the two count-based recipes. In PDZ trans-1, optimized-fit
count overlap is higher than parameter overlap, whereas posterior count
overlap is lower. These comparisons do not rank the sampling recipes
by quality. They show that disagreement with the reported score summary
and uncertainty within an analysis vary separately.

\input{paper/figures/reference_comparisons.tex}

\begin{figure}[p]
\centering
\includegraphics[width=\linewidth]{selection_stability.pdf}
\caption{Candidate inclusion and batch overlap. Panels (a)--(c)
show top-100 inclusion frequencies under parameter resampling for
median-score reference ranks 91--100 (blue) and the ten
highest-frequency outside candidates (orange). Outside candidates
are ordered by inclusion frequency; their row numbers are not
sampled ranks. Panels (d)--(f) show mean shared fractions between
distinct sampled top-$k$ batches, for budgets from 10 to 1,000.
Blue, red, and green denote parameter, optimized-fit count, and
posterior count resampling; the vertical guide marks $k=100$.
Each curve represents one run. PDZ trans-1 and TrpB libB use
Exponential--Exponential, and NucB G3 uses Exponential.
Each example has 16,000 parameter and posterior count samples and
1,000 optimized-fit count samples. Shared fraction measures the
number of common members divided by batch size, not the probability
of exact batch agreement.}
\label{fig:selection}
\end{figure}

\paragraph{The selection budget changes the question.}
The lower row of Fig.~\ref{fig:selection} varies the budget within
each analysis rather than fixing it at 100. NucB G3 parameter samples
share only 21\% of a size-10 batch on average, but 91\% of a size-50
batch and 93\% of a size-100 batch (panel f). This agrees with the
rank intervals in Fig.~\ref{fig:ranks}c: identifying a broader group
can be much more stable than specifying its first few positions.
For this analysis, a ranking-resolution statement without the budget
would omit an important part of the result.

The other examples do not show the same budget response. PDZ trans-1
parameter-sample overlap is 56\% at 10 and 71\% at 1,000.
TrpB libB is nonmonotone over the displayed budgets: overlap falls
from 76\% at 10 to 50\% at 50, then rises to 76\% at 1,000
(Fig.~\ref{fig:selection}d,e). These values are fractions of the
batch, not absolute numbers of common members. Larger budgets admit
a new part of the landscape, and agreement among those additional
candidates need not match agreement among the original leaders.

Although all three examples have higher parameter-sample overlap at
1,000 than at 100, the complete results do not show a universal increase.
After taking the median across models within each landscape,
$\widehat C_{1000}>\widehat C_{100}$ in 13 of 18 landscapes (72\%) for
parameter resampling, 14 of 19 (74\%) for optimized-fit count resampling,
and 13 of 18 (72\%) for posterior count resampling. The same five TrpB
libraries, D--H, decrease under all three recipes
(Table~\ref{tab:landscape-set-sizes}). Larger batches include candidates
at a different cutoff, with different competing scores. This endpoint
comparison does not establish monotonicity between budgets; the curves
show the intermediate values for each running example.

\subsection{Score width alone does not determine batch inclusion}
\label{sec:mechanisms}

Lower input counts are associated with wider score intervals in all
158 leading-candidate analyses. But score width alone does not identify
which candidates will replace members of a batch. A low-count candidate
may remain far below the cutoff, while a closely competing candidate
with a narrower interval may repeatedly cross it. To connect precision
with selection, we examine every
reference top-100 member and challenger together, a group we call
the \emph{decision cohort}. Counts are fixed observed features;
intervals and inclusion frequencies summarize the fitted samples.

This comparison is deliberately restricted to the batch and its
plausible replacements. Most candidates in a landscape may be
irrelevant to a particular selection budget; their precision does
not answer whether the leading list is resolved. Conversely, a
display limited to the reference members would omit the candidates
that replace them. The decision cohort retains both sides of that
comparison without changing the full competitor pool used for ranking.

To compare widths on a stated scale, we divide each candidate's
central 90\% score width by the interquartile range (IQR) of its
run's top-1,000 median scores. The IQR is the difference between
the 75th and 25th percentiles of those scores. It puts uncertainty
on the scale of differences among leading candidates, rather than
equating the assays' numerical units. We use the same denominator
for a candidate's median-score difference from the rank-100 cutoff,
which we call its \emph{score margin}. Positive margins place a
candidate above the cutoff, subject to exact ties; negative margins
place it below.

Figure~\ref{fig:explanation} follows the same candidates across
four panels per row: score width against input count, output count,
and score margin, then inclusion frequency against margin. Filled
blue points are reference members and open orange points are
challengers. Every candidate is retained, without bins, point
subsampling, or pooling across analyses.

The columns separate three aspects of the selection result. Input
counts describe how well the starting composition was observed;
output counts also depend on what happened during selection.
The margin places a candidate relative to the proposed cutoff,
while inclusion frequency records the consequence of all scores
varying jointly. The same count--width association can therefore
coexist with different margin--inclusion patterns. The IQR scaling
allows the plots to be read on a stated relative score scale;
it does not make growth and fluorescence the same physical quantity.

In PDZ trans-1, score widths decline with count support, but span a
narrow range (Fig.~\ref{fig:explanation}a,b). Reference members and
challengers have overlapping widths (panel c), yet candidates with
similar widths have very different inclusion frequencies at different
margins (panel d). There are 319 challengers. The strong count--width
relationship does not make those candidates' selection frequencies
interchangeable: their position relative to the cutoff also matters.

In TrpB libB, the width--count relationship spans a wider range
(Fig.~\ref{fig:explanation}e,f). Some challengers have median scores
well below the cutoff but enter the batch in at least 1\% of samples
(panels g,h). Their broad score distributions make uncertainty in
those low medians relevant to selection. The 416 challengers would
be absent from a display restricted to the reference top 100.
This is not a reason to prefer uncertain candidates; it shows
which alternatives a single median-score list leaves unexamined.

NucB G3 has a different margin--width pattern: many reference members
have wider score intervals than the challengers
(Fig.~\ref{fig:explanation}k). Yet 77 of its 100 reference members
(77\%) have inclusion frequency at least 0.99, and only 28 outside
candidates meet the challenger threshold (panel l). A broad marginal
score interval does not necessarily imply uncertain batch membership.
Selection depends on a candidate's score relative to its competitors
within each joint draw, not on its interval in isolation.

The number of challengers and the number of replacements per sample
are different quantities. TrpB libB has 416 challengers at the 0.01
inspection threshold, but a parameter-sampled top-100 batch contains
36.56 candidates outside its median-score reference on average.
That value is 100 minus the \emph{Reference} overlap of 63.44 in
Table~\ref{tab:reference-comparisons}; it counts all outside candidates,
not only the displayed challengers. NucB G3 has 28 challengers and
an average of 4.64 outside-reference members per sample. A large
collection of possible replacements need not imply that many are
selected together. The threshold determines which alternatives we
inspect, not a new batch-selection rule.

\begin{figure}[p]
\centering
\includegraphics[width=\linewidth]{explaining_resolution.pdf}
\caption{Count support and plausible batch replacements. Rows show
PDZ trans-1 and TrpB libB (both Exponential--Exponential),
and NucB G3 (Exponential), each with 16,000 parameter samples.
Every reference top-100 member (blue, filled) and outside challenger
with inclusion frequency at least 0.01 (orange, open) is retained:
419, 516, and 128 candidates, respectively. Columns show central
90\% score width against input count, output count, and median-score
margin from the top-100 cutoff, then top-100 inclusion frequency
against that margin. Widths and margins are divided by the run's
top-1,000 median-score IQR. Count axes show $\log_{10}(n+1)$; width
axes are logarithmic. Count and width limits are common across
rows. Margin limits are common within a row but differ between
rows. Dotted lines mark zero margin. All candidate points are
shown and no regression lines are fitted. Supplementary
Section~\ref{app:figure-explanation} gives the display conventions.}
\label{fig:explanation}
\end{figure}

Across all runs, within-cohort Spearman correlations between input
counts and normalized score width are negative in 158 of 158
analyses; output-count correlations are negative in 157
(Fig.~\ref{fig:counts}). These are not 158 independent replications:
multiple models and sampling recipes share the same observations.
Output counts also reflect selection, whereas input counts precede
it. Supplementary Section~\ref{app:per-analysis-why} shows the
unaggregated relationships for every run, including Gamma fits.

Descriptive regressions also relate rank width to score width,
local score crowding, their interaction, and reference rank.
Crowding measures how closely spaced neighboring median scores
are. Its association and its interaction with score width vary
across analyses; the interaction is not consistently positive
(Supplementary Section~\ref{app:associations}). Neither those fits
nor the count correlations establish how extra sequencing would
change batch stability. They describe the candidates selected into
the decision cohort under the current analysis.

\section{Discussion}
\label{sec:discussion}

We distinguish disagreement with reported fitness rankings from
uncertainty in ranks and batches within a fitted analysis.

\paragraph{Limitations.}
All uncertainty estimates are conditional on the count models,
score definitions, and retained fits. A model that does not
reproduce the observations can give misleading rank intervals
and inclusion frequencies. Sampler diagnostics were run for all retained
posterior fits, but checking whether a fitted distribution reproduces the experiment
is a separate question. Predictive checks remain incomplete
(Supplementary Section~\ref{app:fit-checks}). Agreement between
sampling recipes is not independent validation, because the
recipes share modeling assumptions. Some results change strongly
with the specification, as illustrated by the PDZ Gamma fits
in Supplementary Section~\ref{sec:sensitivity}. The datasets are
also not a representative sample of all protein assays.

\paragraph{Implications for prediction.}
For machine-learning researchers, these results concern the labels
used to learn and evaluate sequence--fitness relationships. A
reported score is an estimate of an assay-defined quantity, not
an uncertainty-free property of a sequence. A predictor's failure
to reproduce an exact ordering can reflect disagreement over
candidates that the experimental analysis itself does not separate
reliably. Conversely, high correlation over a large cohort does
not guarantee agreement on a smaller batch. These possibilities
motivate evaluating the particular retrieval or ranking decision,
with its candidate set and experimental budget specified.

We have not tested whether label uncertainty changes a comparison
between predictors. A direct next analysis would hold predictions
fixed and evaluate them across sampled score sets, measuring uncertainty
in rank correlation, retrieval, and differences between predictors.
Retraining on uncertain labels would address a separate question.

\paragraph{Batch selection.}
For protein engineers, inclusion frequencies distinguish repeatedly
selected candidates from candidates with plausible replacements
near or below the cutoff. But turnover in identity is not itself
a loss in engineering utility. Alternative variants may have
similar activity, and selection may also value diversity, cost,
or information for a later round. The median-score list is a
reference for measuring turnover, not an optimal rule for every
objective. Assessing the expected activity or information gained
from a batch requires specifying that objective and, where
possible, testing it against follow-up measurements.

\paragraph{Further evidence.}
Two kinds of additional evidence would strengthen the conclusions.
First, predictive checks and suitable held-out measurements could
test whether the fitted models reproduce the experimental
variation relevant to the leading candidates. Second, controlled
read thinning or leave-one-replicate-out fits could test how
changing count support changes batch stability. These require
altered or additional observations; using fewer existing model
draws only assesses Monte Carlo variability.

Under parameter resampling, NucB G3 top-100 batches share about
93 members with each other but only 68 with the reported batch,
whereas TrpB libB batches share about 51 members with each other
(Table~\ref{tab:reference-comparisons}). Disagreement with the reported
list and uncertainty within a fitted analysis are distinct results.

\section*{Data and code availability}
PDZ processed counts and scores are available through GEO accession GSE265816.
TrpB data and analysis code are deposited at
\url{https://doi.org/10.22002/h5rah-5z170}. NucB data and source calculations
are available from \url{https://github.com/google-deepmind/nuclease_design}.
The rank-analysis code consumes prepared sampled-score matrices and aligned
count features; fitting the original models is not required to recreate its
tables and figures. The accompanying repository includes a small pipeline
example, distinct from the real-data analyses reported here.
\note{TODO: Add versioned fitting and analysis repositories and a public
archive for prepared samples, features, and fit diagnostics.}

\section*{Author contributions}
\note{TODO: Confirm contributions with all authors.}

\section*{Competing interests}
\note{TODO: Obtain declarations from all authors.}

\section*{Acknowledgments}
\note{TODO: Confirm funding and acknowledgments.}

\clearpage
\input{paper/references.tex}

\clearpage
\input{paper/materials_methods.tex}

\clearpage
\appendix
\renewcommand{\thesection}{S\arabic{section}}
\renewcommand{\thefigure}{S\arabic{figure}}
\renewcommand{\theHfigure}{supp.\arabic{figure}}
\setcounter{figure}{0}

\input{paper/supplementary_details.tex}
\input{paper/supplementary_analyses.tex}

\end{document}